% Осень26!B13; нумерация внутри столбца без пустых строк.
% Источник: лекция 4, §9 (пространство с мерой)
% Источник: лекция 2, §5 (строка 148)
% Источник: лекция 2, §11 (строка 313)
\definitionticket{10}{Объем, мера, пространство с мерой}
\label{definition:10}

\begin{definition}
Пусть $\mathcal P$ --- полукольцо подмножеств $X$.
Функция $\mu:\mathcal P\to[0,+\infty]$ называется \textbf{объёмом}, если $\mu(\varnothing)=0$ и для любых попарно непересекающихся $A_1,\ldots,A_n\in\mathcal P$, объединение которых $A$ принадлежит $\mathcal P$, выполнено
\[
  A=\bigsqcup_{i=1}^{n}A_i
  \quad\Longrightarrow\quad
  \mu(A)=\sum_{i=1}^{n}\mu(A_i).
\]
\end{definition}

\begin{definition}
Пусть $\mathcal P$ --- полукольцо подмножеств $X$.
Функция $\mu:\mathcal P\to[0,+\infty]$ называется \textbf{мерой}, если $\mu(\varnothing)=0$ и для любых попарно непересекающихся $A_i\in\mathcal P$, $i\in\mathbb N$, объединение которых $A$ принадлежит $\mathcal P$, выполнено
\[
  A=\bigsqcup_{i=1}^{\infty}A_i
  \quad\Longrightarrow\quad
  \mu(A)=\sum_{i=1}^{\infty}\mu(A_i).
\]
\end{definition}

\begin{definition}
\textbf{Пространство с мерой} --- тройка $(X,\mathcal A,\mu)$, где $X$ --- множество, $\mathcal A$ --- $\sigma$-алгебра его подмножеств, $\mu$ --- мера на $\mathcal A$.
\end{definition}
